\chapter{Implementação}
\label{cap:implementacao}

% =====================================================================
% PAPEL DESTE CAPÍTULO (registro de engenharia):
%   Materializa a arquitetura do Cap. 4 em código. Trechos comentados,
%   nível acessível ao leitor professor de Matemática; detalhes extensos
%   vão aos apêndices. NÃO re-explica papéis dos agentes (Cap. 4) nem
%   fundamentos (Caps. 2-3).
%
%   PENDENTE: as Seções 5.6 (custo/desempenho) e 5.7 (exemplos com API
%   real) dependem de execuções com provedor de LLM real e estão
%   marcadas no texto.
% =====================================================================

Este capítulo descreve a implementação concreta da arquitetura apresentada no
Capítulo~\ref{cap:arquitetura}. Antes das escolhas técnicas, cabe explicitar o método que
organizou o desenvolvimento. A pesquisa adota a abordagem de \emph{design science research}
(DSR), apropriada ao mestrado profissional por unir a construção de um artefato à produção de
conhecimento sobre ele, articulada em ciclos: o \textbf{Ciclo 1} construiu o protótipo mínimo
--- um tópico (função quadrática), apenas Gerador e Verificador Simbólico --- para validar a
viabilidade técnica da verificação simbólica em malha fechada com um LLM; o \textbf{Ciclo 2}
completou o sistema com o Crítico Didático, o Orquestrador, o banco curado, a montagem de
listas e a interface do professor, expandindo o recorte para os cinco temas; o \textbf{Ciclo 3},
de avaliação por painel docente, é objeto do Capítulo~\ref{cap:avaliacao}; e o \textbf{Ciclo 4}
corresponde à consolidação do produto educacional e desta dissertação.

\section{Stack técnica}
\label{sec:stack}

O sistema é implementado em \textbf{Python 3.11+}, com quatro dependências centrais:
\textbf{SymPy} para a verificação simbólica (Seção~\ref{sub:sympy}); \textbf{Pydantic} para os
contratos de dados validados por esquema; \textbf{FastAPI} para a interface web do professor;
e o SDK do provedor de LLM escolhido. As dependências de provedor são \emph{opcionais}: o
professor instala apenas a do provedor que usará --- Anthropic, OpenAI ou nenhuma, no caso do
Ollama, que expõe modelos abertos locais por HTTP. Os testes usam \texttt{pytest} e não exigem
chave de API alguma, ponto retomado na Seção~\ref{sec:reprodutibilidade}.

Conforme antecipado na Seção~\ref{sec:decisoes}, não se adota \emph{framework} de agentes: a
orquestração é código Python direto, decisão que privilegia a legibilidade da política de
decisão sobre a conveniência de abstrações prontas.

\section{Estrutura do repositório}
\label{sec:repositorio}

O código vive no diretório \texttt{sistema/} do repositório público da dissertação, organizado
para que cada conceito do Capítulo~\ref{cap:arquitetura} tenha endereço único:

\begin{verbatim}
sistema/
|- dados/bncc_em_matematica.json    # catalogo de habilidades do recorte
|- prompts/gerador.md, critico.md   # prompts versionados (Apendice A)
|- src/questoes/
|  |- especificacao.py              # entrada do professor (Secao 4.2)
|  |- modelos.py                    # contratos entre agentes (Pydantic)
|  |- llm/                          # provedores plugaveis
|  |- agentes/                      # gerador, verificador, critico,
|  |                                #   orquestrador (Secoes 4.3-4.6)
|  |- verificacao/                  # estrategias SymPy por tipo
|  |- banco.py, listas.py           # persistencia e listas (Secao 5.4)
|  |- sincronizacao.py              # espelho na pasta da nuvem (Cap. 7)
|- api/main.py                      # API HTTP da interface (Cap. 7)
|- web/                             # pagina do professor (HTML, CSS, JS)
|- tests/                           # 56 testes automatizados
\end{verbatim}

Dois princípios de organização merecem nota. Primeiro, os \emph{prompts} não são cadeias de
caracteres embutidas no código, mas arquivos Markdown próprios: são artefatos da pesquisa,
citáveis e versionáveis independentemente da lógica. Segundo, o pacote \texttt{verificacao/} é
isolado dos agentes: o núcleo simbólico pode ser testado --- e estendido a novos temas --- sem
tocar em nenhum componente que envolva LLM.

\section{Implementação dos agentes}
\label{sec:impl-agentes}

\subsection{Contratos de dados}

Toda a comunicação entre agentes usa modelos Pydantic, que validam estrutura e tipos no momento
da criação. O contrato central é a formalização verificável (Seção~\ref{sec:ag-gerador}):

\begin{verbatim}
class ExpressaoVerificavel(BaseModel):
    tipo: str            # roteia a estrategia: 'equacao',
                         #   'funcao', 'progressao'
    expressao: str       # ex.: "Eq(2*x**2 - 5*x + 3, 0)"
    incognitas: list[str]
    resposta_esperada: str   # gabarito em sintaxe SymPy
    parametros: dict[str, str]
\end{verbatim}

\subsection{Gerador}

O Gerador monta o pedido ao LLM a partir da especificação --- cada parâmetro vira uma linha de
instrução, com descritores operacionais para dificuldade e natureza --- e valida a resposta
contra o modelo \texttt{Questao}. A extração do JSON é tolerante a desvios comuns de formato
(cercas de código, texto circundante), e qualquer resposta inválida gera erro detectável, nunca
dado corrompido. Nas iterações de revisão, o \emph{feedback} do Orquestrador é anexado ao
pedido com a instrução de corrigir exatamente o que foi apontado, preservando o restante.

\subsection{Verificador Simbólico}

O coração do Verificador é a comparação simbólica exata. Para equações, o conjunto-solução
calculado pelo SymPy é confrontado com o gabarito elemento a elemento:

\begin{verbatim}
def _pertence(valor, colecao) -> bool:
    """Igualdade simbolica exata: simplify(a - b) == 0."""
    return any(sp.simplify(valor - outro) == 0
               for outro in colecao)
\end{verbatim}

\noindent A verificação falha tanto se o gabarito contém solução espúria quanto se omite alguma
--- os dois sentidos da divergência são reportados na justificativa. Um cuidado transversal de
segurança e robustez: as expressões vindas do LLM são interpretadas por \texttt{sympify} com
espaço de nomes \emph{restrito} (apenas símbolos e funções matemáticas autorizados), e qualquer
falha de interpretação degrada para o veredicto \emph{não-verificável} em vez de interromper o
ciclo (Seção~\ref{sec:ag-verificador}).

\subsection{Crítico Didático}

O Crítico envia ao LLM a questão (sem a formalização, irrelevante para o julgamento didático) e
a especificação declarada, sob a rubrica da Seção~\ref{sec:criterios} traduzida em instruções.
A resposta é validada contra o modelo \texttt{ParecerDidatico}, que impõe notas inteiras de 1 a
5 por construção.

\subsection{Orquestrador}

A política de decisão inteira (Seção~\ref{sec:ag-orquestrador}) cabe em um laço legível:

\begin{verbatim}
for numero in range(1, self._max + 1):        # max = 3
    questao = self._gerador.gerar(spec, feedback=feedback)
    verificacao = self._verificador.verificar(questao)

    if verificacao.veredicto == Veredicto.REJEITADO:
        feedback = ...   # veredicto + valor calculado
        continue         # NAO consulta o Critico

    parecer = self._critico.avaliar(questao)
    if not parecer.aprovado:
        feedback = ...   # sugestoes de revisao do Critico
        continue

    return ResultadoCiclo(aprovada=True, ...)
return ResultadoCiclo(aprovada=False, ...)     # descarte
\end{verbatim}

\section{Persistência e banco de questões}
\label{sec:persistencia}

Questões aprovadas são persistidas em \textbf{SQLite} --- formato de arquivo único, sem
servidor, adequado ao uso individual pelo professor. Cada registro guarda a questão completa,
os metadados da especificação como colunas indexáveis (tema, habilidade, nível, dificuldade,
natureza, formato) e o \textbf{histórico integral de iterações} do ciclo que a produziu. Os
metadados alimentam a montagem de listas: o professor filtra o banco por tema e dificuldade,
escolhe a quantidade, e o sistema exporta a lista em Markdown ou LaTeX, em duas versões ---
aluno (sem gabarito) e professor (com gabarito e resoluções).

\section{Reprodutibilidade, logging e testes}
\label{sec:reprodutibilidade}

Três mecanismos respondem ao requisito de reprodutibilidade do projeto. Primeiro, a
\textbf{temperatura} é fixada em 0,3 --- baixa, conforme a discussão da
Seção~\ref{sub:inferencia} --- nos provedores que ainda expõem esse parâmetro. A ressalva não é
detalhe de implementação: as gerações mais recentes de modelos comerciais
\textbf{removeram os parâmetros de amostragem} da interface pública, substituindo o controle
explícito de variabilidade por regulagens internas do próprio modelo. Enviá-los passou a ser
erro de requisição, e o sistema deixa de enviá-los nesses casos. O episódio expõe uma
fragilidade estrutural da dependência de APIs proprietárias: uma decisão metodológica
documentada pode ser invalidada por uma mudança unilateral do fornecedor, sem aviso e sem
recurso --- argumento adicional em favor do caminho com modelos abertos locais, discutido no
Capítulo~\ref{cap:produto}.

Os outros dois mecanismos são independentes do fornecedor e absorvem essa perda. Todo ciclo é
registrado em \emph{log} JSONL com os objetos completos de cada iteração, permitindo
reconstruir qualquer execução --- inclusive quais parâmetros foram efetivamente enviados. E os
\emph{prompts} são versionados junto ao código: qualquer mudança neles é uma mudança rastreável
do sistema.

A estratégia de testes explora a separação de responsabilidades da arquitetura. O núcleo
simbólico é testado diretamente, com casos que incluem os erros típicos de LLM (gabarito
incompleto, troca de sinal) e casos degenerados (PG de razão unitária, extremo de função afim).
A política do Orquestrador é testada com um \textbf{LLM falso} --- um provedor que devolve
respostas pré-programadas em sequência --- que permite exercitar deterministicamente os quatro
caminhos do ciclo: aprovação direta, revisão por rejeição matemática, revisão por reprovação
didática e descarte. O conjunto, de 36 testes, executa em menos de um segundo e não requer
chave de API --- qualquer pessoa que clone o repositório pode validar o sistema integralmente
sem custo.

\section{Custo e desempenho}
\label{sec:custo}

% PENDENTE: esta seção requer execuções com provedores reais.
% Medir: custo médio por questão aprovada (por provedor), tempo médio por
% ciclo, distribuição do número de iterações, taxa de descarte.
\emph{[Seção a completar com dados de execução real: custo médio por questão aprovada em cada
provedor (API comercial \emph{versus} modelo aberto local), tempo médio por ciclo, distribuição
do número de iterações até a aprovação e taxa de descarte por tema. As medições serão feitas
sobre o lote de questões gerado para o banco curado.]}

\section{Exemplos completos de execução}
\label{sec:exemplos}

% PENDENTE: substituir pelos registros JSONL de execuções reais com API,
% reproduzindo os objetos completos (questão, veredictos, pareceres).
\emph{[Seção a completar com três execuções reais registradas pelo sistema: (i) questão de
função quadrática aprovada na primeira iteração; (ii) questão de progressão geométrica rejeitada
pelo Verificador --- gabarito omitindo o primeiro termo --- revisada e aprovada; (iii) questão
não-verificável simbolicamente, aprovada com ressalva e marcada para revisão prioritária do
professor. Os roteiros desses três casos, já validados com o LLM de teste, estão descritos na
Seção~\ref{sec:fluxo}.]}
